\documentclass[10pt]{article}
\usepackage[margin=20mm]{geometry}
\usepackage{amsmath,amssymb,mathtools,fontspec,unicode-math,xurl,hyperref}
\IfFontExistsTF{DejaVu Serif}{\setmainfont{DejaVu Serif}}{\setmainfont{Latin Modern Roman}}
\setmathfont{latinmodern-math.otf}
\setlength{\parindent}{0pt}\setlength{\parskip}{6pt}\sloppy
\title{Technical review of
Towards a Fluid Computer}
\author{AI-assisted technical review}
\date{}
\begin{document}
\maketitle
Abstract: The acquired article combines an exposition of two Euler-flow constructions with a new obstruction to critical compatible metrics. Its reachability definition permits computational difficulty to be placed in an input-dependent computable open target, and the oracle argument in Remark 4.4 does not remove that difficulty. This is a Major definitional and logical issue, without refuting the stronger step-by-step simulations actually used in the constructions. The critical-metric argument has a local geometric repair. A suspension construction answers the stated Hodge-viscosity version of Open problem 4 when the compact manifold and metric may be chosen freely. Several local formula and coordinate corrections are also identified.\par
\section*{Source and review scope}
The source is Robert Cardona, Eva Miranda, and Daniel Peralta-Salas, Towards a Fluid Computer, Foundations of Computational Mathematics 25 (2025), 1921–1937, DOI 10.1007/s10208-025-09699-6 (https://doi.org/10.1007/s10208-025-09699-6). The acquired publisher version of record was published online on 13 March 2025 and has 17 PDF pages. References to pages below use the printed pagination.\par
Although much of the article is expository, Proposition 4.3 is presented as a new consequence, and the paper remains within the assigned research-review scope. The constructions imported from the article’s references [7] and [8] are distinguished from arguments supplied here. Their complete external proofs, and the classification results cited as [2] and [18], were not independently retrieved. Missing details explicitly assigned to those works are not by themselves treated as defects.\par
\section*{P1. The computable-target definition does not isolate computation by the dynamics (Major)}
Definition 2.2, p. 1924, permits a computable point \(p_c\) and a computable open set \(U_c\) depending on the input \(c\), with halting equivalent to the positive orbit meeting \(U_c\). The surrounding paragraph describes computability of an open set through arbitrary-precision approximation, without restricting its geometric presentation. Under the usual computable-real endpoint or effective-open interpretation, even the identity map meets this definition.\par
Fix the universal machine in the definition. Number its steps from one, and put \[r_c=
 \begin{cases}
  2^{-t},&\textstyle\text{if the machine first halts on }c\text{ at step }t,\\
  0,&\textstyle\text{if it never halts}.
 \end{cases}\] The real number \(r_c\) is uniformly computable. To approximate it with error at most \(2^{-N}\), simulate \(N\) steps. If halting occurs, return \(2^{-t}\); otherwise return zero. A later halting time produces an error at most \(2^{-(N+1)}\). This approximation procedure does not decide whether \(r_c\) is zero.\par
In the unit disk take the identity map, the fixed rational point \(p_c=(0,0)\), and the nonempty open rectangle \[U_c=(-r_c,1/4)\times(-1/4,1/4).\] Since \(0\leq r_c\leq1/2\), this rectangle lies in the disk. It has uniformly computable endpoints and is uniformly effectively open; its endpoint description can be approximated to every prescribed accuracy. Nevertheless, \[\{\operatorname{Id}^n(p_c):n>0\}\cap U_c\ne\varnothing
 \quad\Longleftrightarrow\quad r_c>0
 \quad\Longleftrightarrow\quad T_U(c)\text{ halts}.\] The orbit does nothing. All undecidability is carried by a boundary-membership question for the target.\par
The same problem occurs directly in the contact setting. Let \(R=\partial_{\phi_1}+\partial_{\phi_2}\) be the standard Hopf field on the round \(S^3\), take the rational point \((z_1,z_2)=(3/5,4/5)\), and set \[U_c=\{\,|z_1|^2>9/25-r_c\,\}.\] The quantity \(|z_1|^2\) is constant on the orbit, so reachability is again equivalent to \(r_c>0\). The Hopf field is Killing for the standard compatible metric, and therefore its Chern–Hamilton energy is zero. Thus the weak target-based definition cannot support a general exclusion of these flows.\par
Remark 4.4, p. 1930, assumes an oracle deciding whether two computable sets intersect, then asserts that the relevant halting problem remains undecidable. Applied to \(\{p_c\}\) and \(U_c\), that oracle already decides ordinary halting in the example above. Undecidability of halting relative to the stronger oracle is a different assertion; no simulation of all oracle machines by the fixed constructed system is supplied. The oracle comparison consequently does not establish the claimed exclusion.\par
This finding does not refute the actual injective, configuration-by-configuration simulations described in Sections 4 and 5. Those have substantially more structure than Definition 2.2. A repair should specify a uniform effective simulation and a restricted decoding or geometric target representation, and verify that the constructions satisfy it. Merely assuming a set-intersection oracle, or merely making an otherwise unrestricted effective-open target independent of the input, does not resolve the issue. If a stronger finite target representation was intended, it must be stated; the displayed counterexample addresses the finite-precision meaning explicitly suggested in the article.\par
\section*{P2. Proposition 4.3 needs a section-independent argument (Minor)}
The proof on pp. 1929–1930 uses the cited classifications to obtain orbital conjugacy to \[R_{l,k}=l\partial_{\phi_1}+k\partial_{\phi_2},\qquad l,k>0.\] On the standard section \(\{\phi_1=0\}\), its return angle is correctly given by \(2\pi k/l\). The proof then infers that this angle vanishes modulo \(2\pi\) because the constructed return map on a transverse disk is the identity near its boundary. Those need not be the same section.\par
For example, \(R_{2,1}\) has return rotation by \(\pi\) on \(\{\phi_1=0\}\), whereas its return map on \(\{\phi_2=0\}\) is the identity: a return takes time \(2\pi\) and advances \(\phi_1\) by \(4\pi\). Thus the existence of one identity return map does not determine the rotation number of the other.\par
The conclusion for the selected universal construction can be repaired as follows. The identity collar of its disk return map suspends to an open region of periodic flow orbits. An irrational weighted Hopf flow has only two exceptional closed circles, so orbital conjugacy forces \(l/k\) to be rational. Every orbit of the weighted Hopf flow is then closed, and the same is true for the orbitally conjugate field.\par
On the other hand, a universal reversible machine has an initial nonhalting run with no repeated configuration. Otherwise, simulation until either halting or a repeated finite configuration would decide halting on every initial input. The injective Cantor encoding and step-by-step conjugacy therefore give an aperiodic return orbit. A closed flow orbit cannot supply this aperiodic orbit of the constructed transverse return map. This is the required contradiction.\par
This repair uses the actual simulation and the periodic collar, without the weak reachability definition or Remark 4.4. It is conditional on the two quoted classification results, which remain imported dependencies. Proposition 4.3 should explicitly refer to the selected universal construction of Corollary 4.2; Theorem 4.1 also realizes arbitrary disk maps, including the identity, so its full class of Reeb fields cannot satisfy that exclusion. The claim concerns critical compatible metrics, not every metric for which an arbitrary field happens to be Beltrami.\par
\section*{P3. The stated Hodge-viscosity open problem has an affirmative suspension construction (Minor)}
Open problem 4, p. 1935, asks for a Turing-complete steady Navier–Stokes flow on some compact Riemannian three-manifold. Remark 6.1 specifies the Hodge Laplacian as the viscous operator. With this freedom in the manifold and metric, the area-preserving simulation already used in the article supplies such a flow.\par
Let \(\varphi\) be the simulated area-preserving disk diffeomorphism from Proposition 4.1, equal to the identity in a boundary collar. Extend it by the identity to a closed oriented surface \(\Sigma=S^2\) with area form \(\omega\) agreeing with the disk area. Form the mapping torus \[M=(\Sigma\times[0,1])/\bigl((x,1)\sim(\varphi(x),0)\bigr).\] The vector field \(X=\partial_t\), the one-form \(\alpha=dt\), and the volume form \(\mu=dt\wedge\omega\) descend smoothly, because \(\varphi\) preserves \(\omega\).\par
Here is an explicit compatible metric construction. Choose a surface metric \(h_0\) with area form \(\omega\), and choose a smooth function \(\beta:[0,1]\to[0,1]\) that is zero near \(t=0\) and one near \(t=1\). Put \[k_t=(1-\beta(t))h_0+\beta(t)\varphi^*h_0,\qquad
 \operatorname{vol}_{k_t}=f_t\omega,\qquad h_t=f_t^{-1}k_t.\] The convex combination is positive definite. Since the surface has dimension two, \(\operatorname{vol}_{h_t}=\omega\). Also \(h_1=\varphi^*h_0\), and the family is constant near its endpoints. Consequently \[g=dt^2+h_t\] descends to a smooth metric on the closed compact mapping torus.\par
For this metric, \[X^\flat=\alpha,\quad |X|_g=1,\quad
 \operatorname{vol}_g=\mu,\quad
 \operatorname{div}_g X=0,\quad
 d\alpha=0,\quad \delta_g\alpha=0.\] Thus \(\Delta_H\alpha=0\), and the covariant identity \[(\nabla_X X)^\flat=\iota_Xd\alpha+\tfrac12d|X|^2\] gives \(\nabla_X X=0\). Constant pressure and \(X\) solve the unforced steady Navier–Stokes equations for every viscosity, since the Hodge viscous term vanishes. Its sign convention is immaterial here. The kinetic energy is \(\operatorname{Vol}_g(M)/2\), hence finite.\par
The section \(t=1/2\) has return map \(\varphi\). A computable initial point on that section and a halting target lifted to a fixed product strip, for example \(3/4<t<7/8\), reproduce the discrete orbit’s target visits. This transfers the actual step-by-step simulation, not merely the weak boundary-encoding device of P1.\par
The construction answers the precise arbitrary-metric, Hodge-viscosity formulation stated in the paper. It does not produce a solution for a prescribed round or flat metric, prove robustness under perturbations, or apply automatically to a different geometric viscous operator. Those restrictions could lead to a substantive revised question. The observation is classified as a bounded correction to the open-problem formulation, not a failure of the principal Euler constructions.\par
\section*{P4. Local Euler and encoding formulas need correction (Minor)}
On p. 1925 the Bernoulli function is printed as \(B=p+\|X\|/2\). The correct formula is \[B=p+\tfrac12|X|^2.\] Indeed \((\nabla_X X)^\flat=\iota_Xd(X^\flat)+\tfrac12d|X|^2\). An exact symbolic check with \(X=(-y,x,0)\) and \(p=(x^2+y^2)/2\) verifies stationary incompressible Euler. At \((1,0,0)\), the components of \(\iota_Xd(X^\flat)\) are \((-2,0,0)\), whereas \(-d(p+|X|/2)\) gives \((-3/2,0,0)\). The squared expression agrees identically. This is a missing exponent, not an obstruction to the corrected Beltrami correspondence.\par
In the planar encoding, pp. 1932–1933, the one-dimensional interval should be translated by \(2j\): \(I_{(q,t)}^j=I_{(q,t)}+2j\), rather than by the displayed two-dimensional vector \((2j,2j)\). The initial point should likewise be \[p_{c_i}=(2i+\phi(c_i),0),\] to belong to the designated strip \([2i,2i+1]\times[0,\infty)\). The printed \(i+\phi(c_i)\) fails already for \(i=1\). These corrections make the encoding agree with the subsequent curves and figures.\par
\section*{Positive checks and limits of the review}
The finite-tape encoding in Section 4 gives rational ternary coordinates. The four affine branches in Fig. 1 are consistent with its specified \(F,G\) rules: on the indicated blocks they are \[(3x,y/3),\quad(3x-2,y/3-2/9),\quad
 (x,y+2/3),\quad(x,y+2/3).\] Each has determinant one, and the image strips and corner blocks agree with the actual figure. The bijective generalized-shift premise and reversibility are essential when extending to a disk diffeomorphism.\par
Cartan’s formula verifies preservation of \(\alpha\) and \(\alpha\wedge d\alpha\) by a Reeb field. A positive rescaling \(hR\) preserves \((\alpha\wedge d\alpha)/h\). Together with an adapted metric, the corrected Bernoulli identity yields stationary Euler. This explains the flexibility of the non-prescribed-metric construction; it does not yield a fixed Euclidean metric result.\par
Section 5 is explicitly a sketch of a different imported construction. The reciprocal prime-power encoding is injective; its intervals can be separated with radii of order the square of their centers. The stated weak robustness uses an error that decays sufficiently rapidly at infinity, and is not uniform perturbation robustness on a compact phase space. The absence of a uniform lower bound on the decoding intervals is explicitly acknowledged.\par
For the planar Cauchy datum \((F_x,F_y,0)\), the normal curl equation requires \(\partial_xF_y-\partial_yF_x=0\), as stated. Theorem 5.2 requires extension of \(F\) to an entire function on \(\mathbb C^2\), not merely real analyticity. This distinction is retained when using the global extension. Tangency makes the plane invariant, so its simulated trajectories remain trajectories of the Beltrami field. The full entire-approximation and global Cauchy–Kovalevskaya results are cited dependencies, not newly certified external theorems. The Euclidean construction’s infinite-energy limitation is acknowledged by the article and is not reported as a hidden defect.\par
The assessment is Major revision because the computational definition and its exclusion argument require substantive correction. The specific fluid constructions are not shown false. Proposition 4.3 has the conditional repair above, and the additional findings are local or concern the precise scope of an open question. No claim is made to resolve the paper’s fixed-metric Euler, entropy, or celestial-mechanics questions.\par
\section*{Scope and references}
No substantive earlier report for M25-0449 was recoverable in the bounded preserved inventory; inaccessible conversation-only material is recorded as unavailable. The acquired source PDF, manifest, and historical extraction were reused without relabeling their execution. The current review, actual source-page views, symbolic local identity check, and report build are new. Analytic computability and geometric arguments are recorded as arguments, not as numerical simulations.\par
The exact source, scripts, outputs, helper roles, build inputs and retained dependencies, and actual current report-page inspections are bound in the accompanying verification files. No author code was executed, and absent optional code is not a mathematical defect. Scientific disposition remains pending; no human review or independent acceptance is recorded.\par
1 R. Cardona, E. Miranda, and D. Peralta-Salas, Towards a Fluid Computer, Found. Comput. Math. 25 (2025), 1921–1937. Published online 13 March 2025. doi:10.1007/s10208-025-09699-6 (https://doi.org/10.1007/s10208-025-09699-6).\par
\end{document}
